\noindent \resheading{Research Experience}
\begin{itemize}

\item \cvdated{CT-less PET Imaging using the J-PET Scanner \normalfont{(PhD Topic)}}{Oct 2024--Present} \\
\textit{PhD Supervisors}: \textbf{Prof. Paweł Moskal \& Dr. Sushil Sharma}, Jagiellonian University
    \begin{itemize}
        \item Developing CT-less PET reconstruction using \textbf{scattered and annihilation data} from the plastic-scintillator J-PET scanner.
        \item Exploring deep learning for attenuation-map denoising and mutual-information-based techniques for PET image alignment.
        \item Targeting improved lesion localization and image quality without anatomical priors from CT or MRI.
    \end{itemize}

\item \cvdated{Using Scattered Data in PET Scan Imaging \normalfont{(MSP-1)}}{July 2023--Nov 2024} \\
\textit{Project Supervisor}: \textbf{Prof. Pragya Das}, Dept. of Physics, IIT Bombay
    \begin{itemize}
    
        \item Validated Dr. S. Ghosh's approach for tumor localization using energy-based scattered data with \textbf{MATLAB} GUI tools.
        
        \item Developed a custom 2D TOF-based model and collaborated with peers to test its accuracy using artificial phantom data.
        
        \item Used Desmos-based visualizations to evaluate the mathematical model qualitatively.
        
    \end{itemize}

\item \cvdated{Designing a 3D Model for Scattered Data in PET Imaging \normalfont{(MSP-2)}}{Jan 2024--Apr 2024} \\
\textit{Project Supervisor}: \textbf{Prof. Pragya Das}, Dept. of Physics, IIT Bombay
    \begin{itemize}
    
        \item Applied TOF method to estimate annihilation points, enhancing geometric precision in 3D PET reconstruction.
        
        \item Used \texttt{meshgrid} and \texttt{fsolve} to solve sphere-line intersection problems in 3D space.
        
        \item Proposed model improvements, identified limitations in data, and outlined future research directions in PET imaging.
        
    \end{itemize}

\end{itemize}